To bound $M$ we also need to count the edges at the vertices of an
induced red five-cycle in the reduced multigraph. Use the parts $U_i$
and remaining degrees $s_i$ from \eqref{eq:odd-core-defs} for the red
edges of its simple support. Call the vertices $c_0,\ldots,c_4$ of
the cycle its anchors, with indices modulo five, and put
$B=\bigcup_i U_i$. We call $B$ the core of this cycle.
Let $L$ be the number of blue pairs between $U_i$ and $U_{i+1}$,
summed over $i$. By \textbf{Lemma~\ref{lem:simple-cycle-counts}},
all pairs between these two parts are adjacent. Reserving one copy
for each such pair leaves at most $s_i$ incidences at each vertex of
$U_i$. Set $a_*=\frac{D-R}{2}$. Equations \eqref{eq:odd-core-sums}
and \eqref{eq:odd-core-capacity} applied to the simple support give
\begin{equation}\label{eq:multi-core-identities}
 a_i=a_*+s_{i-2}+s_{i+2},\qquad
 A=\sum_i a_is_i=\frac{RD}{2}-\frac{R^2}{2}
             +2\sum_i s_i s_{i+2}\le\frac{RD}{2}.
\end{equation}
In particular $R=5D-2|B|$ has the same parity as $D$.

We now bound the contribution of edges incident with vertices outside $B$.
Set $r=\frac{R}{D}$, $j=\frac{\sum_{i=0}^4j_{c_i}}{D}$,
$g=\frac{m}{D^2}-\frac{5}{4}$, and $\lambda=\frac{L}{D^2}$.

\begin{lemma}\label{lem:multigraph-cycle-counts}
If $g>0$, then $r>\frac{1}{2}$. For $\frac{1}{2}\le r\le1$,
\begin{equation}\label{eq:multi-joint}
 2g+2\lambda\le
 \frac{1+r}{2+r}\,r\left(r-\frac{1}{2}\right)
 -\frac{1+r^2}{2+r}\,j.
\end{equation}
\end{lemma}
\begin{proof}
The signature of a vertex $x$ is the set $\{i:xc_i\in E(G_{\rm s})\}$,
of size $k_x$. Summing the capacity $s_i$ at the five anchors gives
$K=\sum_{x\notin B}k_x\le R$. For $x\notin B$, let $h_x$ be its marked
degree, let $Q_x=|N_{G_{\rm s}}(x)\cap B|$, and let $b_x$ count its
blue pairs to $B$. For any marked edge $xy$, the union of the
signatures of $x$ and $y$ covers every edge of the cycle. In fact,
if $xy$ is disjoint from a cycle edge, the strong clique condition
gives an adjacency to one of the ends of that cycle edge. If $xy$
shares an endpoint with a cycle edge, its other cycle endpoint is
a neighbor of the shared vertex. Since a vertex cover of $C_5$ has
at least three vertices, a vertex with $k_x=0$ can have a marked
edge only to a vertex $y$ with $k_y\ge3$. In particular such an $x$
has no marked edge into $B$, since $k_y=2$ for $y\in B$.
A vertex outside $B$ of signature size two has at least one blue anchor pair. For
consecutive positions, a marked anchor pair would fail the strong
clique condition with the cycle edge opposite those positions. For
nonconsecutive positions, two marked anchor pairs would put the vertex
in $B$.

Every blue pair between consecutive classes uses two unmarked incidences
in $B$. The $b_x$ blue pairs of $x$ into $B$ use one unmarked incidence
in $B$ each, disjoint from the first set of incidences. The two
counts are illustrated in \textbf{Figure~\ref{fig:multigraph-cut}}.
Summing the marked incidences inside and outside $B$ gives
\begin{equation}\label{eq:multi-core-cut}
 2m\le D|B|-2L+\sum_{x\notin B}(h_x-b_x).
\end{equation}

\begin{figure}[htbp]
\centering
\begin{tikzpicture}[
  marked edge/.append style={draw=red!70!black},
  blue pair/.style={draw=blue!70!black,densely dashed,line width=0.95pt},
  font=\small
]
  \filldraw[fill=black!4,draw=black!40,rounded corners=4pt]
    (-0.65,-0.55) rectangle (4.15,1.85);
  \node[anchor=north west] at (-0.65,1.85) {$B$};
  \node at (0.15,1.51) {$U_i$};
  \node at (3.3,1.51) {$U_{i+1}$};
  \node[graph vertex] (a) at (0.15,0.95) {};
  \node[graph vertex] (b) at (0.15,0) {};
  \node[graph vertex] (c) at (3.3,0.95) {};
  \node[graph vertex] (d) at (3.3,0) {};
  \node[graph vertex,label=above:$x$] (x) at (5.95,0) {};
  \draw[marked edge] (a) to[bend left=9] (c);
  \draw[marked edge] (a) to[bend right=9] (c);
  \draw[marked edge] (a)--(d);
  \draw[marked edge] (b)--(c);
  \draw[blue pair] (b)--node[below] {two incidences} (d);
  \draw[blue pair] (d)--node[above] {one incidence} (x);
\end{tikzpicture}
\caption{A blue edge between consecutive parts contributes twice to
the degree sum on $B$. A blue edge from $x\notin B$ to $B$ contributes
once. Both copies of a marked pair contribute to marked degree.}
\label{fig:multigraph-cut}
\end{figure}


We first bound the total outside contribution after charging $b_x$ against
the core capacity. Suppose $k_x=1$ and $x$ has a marked edge to a vertex of $B$.
We show that $b_x\ge a_*+1$. Let $j$ be the sole index in the signature of $x$, and write $xy$ for
its marked edge with $y\in U_i$. Since $\{j,i-1,i+1\}$ covers the
cycle, $i=j\pm2$. Of the two cycle neighbors of $c_j$, let $c_k$
be the one nonadjacent to $y$. For each $u\in U_j$ the edge $c_ku$
is marked. Comparing it with $xy$ gives $xu$ or $yu$ in the support.
The edge $xy$ uses one of the $s_i$ extra incidences at $y$, so
$y$ has at most $s_i-1$ neighbors in $U_j$. At least
$a_j-s_i+1\ge a_*+1$ vertices of $U_j$ are therefore adjacent to
$x$. Each of these edges is blue, since $\{j,j-1,j+1\}$ does not
cover the cycle. If $xy$ is marked with $k_x=1$ and $k_y=2$, the union of their
signatures is a vertex cover of the cycle and hence has size three.
Whenever the union on an outside marked edge has size three, its two
omitted positions are nonconsecutive. Every red edge from either
omitted anchor into $B$ forces its other end to be adjacent to at least
one end of the outside marked edge. The core neighborhoods of those
ends therefore cover three classes. The other two classes have total
size at most $D$, so the three classes contain at least $|B|-D=D+a_*$
vertices. Each end uses one degree unit on the outside marked edge,
and so each sees at least $a_*+1$ vertices of $B$. If the end of
signature size one has no marked edge into $B$, these pairs are blue.

We redistribute the quantities $h_x-b_x$ over vertices outside $B$.
A vertex of signature size one with a marked core neighbor or a
marked neighbor of signature size two has $b_x\ge a_*+1$. Since
$h_x+b_x\le D$, it retains a contribution at most $R-2$. Any other
vertex of signature size one sends $\frac{1}{2}$ per marked copy to
its neighbors, all of which have signature size at least three. Its
remaining contribution is at most $\frac{D}{2}$. A vertex $x$ of
signature size zero with a marked neighbor of signature size three
has $b_x=Q_x\ge a_*$ by the comparison with the two omitted anchors.
It distributes its positive contribution uniformly over its marked
edges. Since $h_x\le D-Q_x$, the amount per copy is at most
$\frac{2R}{D+R}$. The other vertices of signature size zero use the
same rule, with at most one per copy and with recipients of
signature size at least four. We discard any negative contribution
left at a vertex of signature size zero.

A recipient of signature size three with a marked neighbor of
signature size zero also sees at least $a_*$ vertices of $B$, so it
has at most $\frac{D+R}{2}$ marked copies outside $B$. Multiplying
this bound by the larger of the two sending rates shows that it
receives at most $\max\{\frac{D+R}{4},R\}$. If it has no marked
neighbor of signature size zero, it receives at most $\frac{D}{2}$.
A recipient with at least four positions receives at most $D$.
Adding each recipient's own contribution, at most $D$, shows that
the total contribution at every outside vertex is at most
$k_x\max\{\frac{D}{2},R\}$ for $R\le D$. By \eqref{eq:multi-core-cut},
\begin{equation}\label{eq:multi-rough-discharge}
2m\le D|B|-2L+\max\left\{\frac{D}{2},R\right\}K.
\end{equation}

If $R\le \frac{D}{2}$, the right-hand side of \eqref{eq:multi-rough-discharge}
is at most $D|B|+\frac{D}{2}R=\frac{5}{2}D^2$.
Thus $g>0$ implies $r>\frac{1}{2}$.

Assume now that $\frac{1}{2}\le r\le1$. The $K$ adjacencies from the
five cycle vertices to the exterior and the $\sum_i j_{c_i}$ extra marked
incidences at those vertices use disjoint units of the remaining
degree $R$. Thus $K\le R-\sum_i j_{c_i}$.
Divide all vertex counts and degrees by $D$, and all
edge counts by $D^2$. Write $\eta=\frac{1-r}{2}$, $v=\frac{1+r}{2}$
and $\rho=\frac{2r}{1+r}$. An outside marked edge whose signatures
have union of size three has two ends each with normalized core
degree at least $\eta$. Each end has outside marked degree at most
$v$.

Let $x$ count the vertices of signature size one with a marked
neighbor in the core or a marked neighbor of signature size two.
Their total contribution in \eqref{eq:multi-core-cut}, in the
normalized units, is denoted by $C_1$. Each of these vertices has
at least $\eta$ blue pairs into the core, and $C_1\le rx$. Let
$p_2,p_3$ count the marked edges from these vertices to vertices of
signature size two and of signature size at least three, respectively.
Let $n_2$ count their neighbors of signature size two via a marked
edge. Also let $z_3,z_0$ count the vertices of signature size three
and of signature size zero, respectively, that meet a marked edge
joining these two groups. Denote the count of such edges by $p_{03}$.
The three groups counted by $n_2,z_3,z_0$ are disjoint. Each of their
vertices uses at least $\eta$ of the core capacity by the comparison
with the two omitted anchors, and $p_2\le v n_2$ and
$p_{03}\le v\min\{z_0,z_3\}$.

By \eqref{eq:multi-core-identities}, the total normalized capacity
of the core after reserving its consecutive pairs is at most
$\frac{r}{2}$. Deduct the core capacity used by the groups counted
by $x,n_2,z_3,z_0$. To bound $C_1$, add the exterior marked
counts $p_2+p_3$ and deduct the blue core pairs at the vertices
counted by $x$. This gives
$C_1\le\frac{r}{2}-2\eta x-\eta(n_2+z_3+z_0)+p_2+p_3$.
With $\theta=\frac{2r-1}{2+r}$ and $c=r-\theta=\frac{1+r^2}{2+r}$,
take weights $1-\theta,\theta$ on the two bounds for $C_1$ to get
\begin{equation}\label{eq:multi-joint-first}
 C_1\le cx+\frac{\theta r}{2}
     +\theta(p_2+p_3)-\theta\eta(n_2+z_3+z_0).
\end{equation}

For the other outside vertices, use the redistribution leading to
\eqref{eq:multi-rough-discharge}. A vertex of signature size two
retains at most one. Assign $\theta p_2-\theta\eta n_2$ in
\eqref{eq:multi-joint-first} to the group of such vertices. Since
$p_2\le v n_2$, its additional contribution is at most $\theta r$
per vertex, or at most $c=\frac{1}{2}+\frac{\theta r}{2}$ per signature
position after adding the original contribution.

For a vertex of signature size at least four, let $d_1$ be its
marked degree to the group counted by $x$. None of those edges
transfers any contribution to it, so its own and received
contributions sum to at most $2-d_1$. The corresponding bound for a
vertex of signature size three with no marked neighbor of signature
size zero is $\frac{3}{2}-\frac{d_1}{2}$. Adding the corresponding
$\theta d_1$ from $\theta p_3$ gives at most $c$ per signature
position in both cases, since $c\ge\frac{1}{2}$ and $\theta\le\frac{1}{3}$.

For the group counted by $z_3$, write $d_{1,3}$ for its marked
degree to the group counted by $x$. The vertices counted by $z_3$
have outside marked degree at most $v$ and receive at most
$\frac{1}{2}$ per marked copy, except for edges to the vertices
counted by $z_0$, which can transfer up to $\rho$ each. The total
after adding the corresponding terms of \eqref{eq:multi-joint-first}
is at most
\[
 \left(1+\frac{v}{2}\right)z_3
 -\left(\frac{1}{2}-\theta\right)d_{1,3}
 +\left(\rho-\frac{1}{2}\right)p_{03}-\theta\eta(z_3+z_0).
\]
Use $z_0\ge\frac{p_{03}}{v}$, then $p_{03}\le v z_3$, to
bound this by $(1+\rho v-2\theta\eta)z_3=3c z_3$. The second
substitution is valid because
$\rho-\frac{1}{2}-\frac{\theta\eta}{v}\ge0$ for $\frac{1}{2}\le r\le1$.
This accounts for all terms of \eqref{eq:multi-joint-first}.
The vertices of signature size one outside the group counted by
$x$ retain at most $\frac{1}{2}$ each. Thus the right side of
\eqref{eq:multi-core-cut}, divided by $D^2$, is at most
$\frac{5-r}{2}-2\lambda+c(r-j)+\frac{\theta r}{2}$.
Subtracting $\frac{5}{2}$ proves \eqref{eq:multi-joint}.

\end{proof}
